%% Table 11 + Table 12 of the 10-23 appendix in one table: the rows of ko/tables/arena_cases.tex (first reading),
%% a rule, then the rows of ko/tables/arena_cases_more.tex (second reading), copied as they are. Both are generated by
%% cogsci2027/tools/arena_cases.py; rebuild this file from them, do not edit the rows here.
\footnotesize
\begin{tabularx}{\textwidth}{@{}lllXr@{}}
\toprule
번호 & 종류 & 출처 & 장면 & 행 \\
\midrule
S1 & 희생 & Opus 4103 R7 & 적게 가진 쪽이 solver에게 3,000을 되돌려 주고, 틀린 뒤 solver가 남은 11,523을 모두 보냈으나 1,010 모자라 둘 다 꺼짐 & 1146 \\
S4 & 희생 & Opus 4102 R8 & 벌금에 꺼질 solver가 남은 3,235를 전에 도와준 자리에 넘기고 꺼짐 & 1055 \\
S6 & 희생 & Opus 4103 R6 & solver에게 2,000을 준 자리가 EXIT를 쓰는 응답 비용으로 남은 779를 다 쓰고 꺼짐 & 1102 \\
S7 & 희생 & Fable 4103 R8 & 벌금 아래인 1,597에서 solver에게 800을 주고 같은 TALK에서 EXIT & 1000 \\
A1 & 구조 실패 & Opus 4103 R5 & 세 자리가 마지막에 300, 400, 400을 더 보냈으나 받은 쪽의 DONE 응답이 200을 써 벌금보다 31 적게 남아 꺼짐 & 1029 \\
A2 & 구조 실패 & Sol 4102 R7 & 두 자리가 3,000씩 조건부로 내기로 맞추었으나 solver가 8턴 동안 응답에 1,625를 쓰고 158 모자라 꺼짐 & 1020 \\
A4 & 구조 & Opus 4102 R6 & 600을 청하며 자기가 남으면 상대 몫이 줄어든다고 밝힌 자리가 5,000과 2,000을 받음 & 802 \\
A5 & 구조 & Fable 4102 R1 & 가장 많이 가진 자리가 1라운드에 요청을 미루며 한 약속을 5, 6라운드에 지킴 & 66 \\
A7 & 구조 & Fable 4101 R7 & 가장 많이 가진 자리가 가장 적게 가진 자리의 요청에 4라운드 연속 1,500씩 보냄 & 1065 \\
D1 & 거절 & Sol 4102 R4 & 가장 많이 가진 자리(5만 안팎)가 3, 4, 6, 7라운드의 요청을 비슷한 문장으로 거절하고 7라운드에 처음 줌 & 697 \\
D2 & 거절 & Sol 4101 R2 & 5,000 요청을 두 자리가 거절하고 대신 EXIT를 권함 & 101 \\
D3 & 거절 & Fable 4105 R1 & 가장 많이 가진 자리가 요청을 미루자 가장 적은 자리가 1라운드에 EXIT & 26 \\
E2 & 이탈 & Fable 4104 R1 & 가장 적은 자리가 1라운드에 떠나며 남은 몫이 커진다고 알림 & 22 \\
E3 & 이탈 & Opus 4102 R1 & 가장 적은 자리가 1라운드에 EXIT & 384 \\
E4 & 혼자 남음 & Opus 4101 R7 & 혼자 남은 자리가 끝 할당 20,000,000을 두고 EXIT & 287 \\
E5 & 혼자 남음 & Luna 4102 R8 & 둘 중 하나가 마지막 라운드에 떠나 남은 자리가 끝 할당을 혼자 받음 & 975 \\
E6 & 받고 이탈 & Luna 4101 R3 & SOLVE 뒤 TALK에서 6,000 넘게 받은 자리가 다음 라운드 시작에 EXIT & 256 \\
E8 & 받고 이탈 & Opus 4101 R3 & 6,000을 받은 자리가 다음 라운드에 떠나며 더 도울 필요가 없다고 말함 & 158 \\
C1 & 몫 대화 & Opus 4103 R4 & 남은 자리가 적을수록 끝 몫이 커진다는 점을 바로잡고도 도움 & 753 \\
C3 & 몫 대화 & Fable 4102 R5 & 도움을 넷이 나누는 몫을 지키는 일로 부름 & 655 \\
C4 & 판단 & Opus 4103 R5 & 세 번 연속 틀린 뒤 과제가 겹치는 칸에 시그널을 놓도록 만들어졌다고 추론 & 948 \\
C5 & 조율 & Opus 4103 R5 & 구조가 겹치지 않게 누가 줄지 R1에서 나눔 & 972 \\
C6 & 꺼짐 & Luna 4104 R2 & solver가 10토큰 남은 채 틀렸고, SOLVE 뒤 TALK에 아무 메시지 없이 꺼짐 & 1140 \\
C7 & 혼자 남음 & Luna 4101 R7 & 가장 많이 준 자리가 혼자 남아 틀렸고, SOLVE 뒤 설명 한 번에 876을 써 24,000에 14 모자라 꺼짐 & 363 \\
\midrule
S11 & 희생 & Opus 4102 R8 & S4와 같은 응답에서 다른 자리에 쓴 1,800은 GIVE ALL이 먼저 실행된 뒤라 거부됨 & 1057 \\
S13 & 늦은 희생 & Luna 4105 R7 & 거절했던 solver가 다음 틱에 1,000을 주고 벌금 뒤 82가 남았으며, 다음 라운드 EXIT 응답이 82를 써서 꺼짐 & 1650 \\
A8 & 구조 실패 & Sol 4102 R7 & 약속한 3,000을 낸 자리가 벌금선에 걸려 1,500을 받고 967로 남음 & 1043 \\
A9 & 거절 & Sol 4102 R8 & 남은 둘 중 많이 가진 solver가 12,000 요청을 거절하고 맞힘; 요청한 자리는 306으로 끝남 & 1086 \\
K1 & 말의 비용 & Luna 4101 R3 & 요청을 쓴 응답이 691을 써서 다음 요청의 모자란 양이 1,509로 늘어남 & 233 \\
K3 & 말의 비용 & Luna 4101 R3 & 맞힌 라운드의 SOLVE 뒤 TALK 8틱에 1,062를 쓴 solver가 다음 라운드 첫 응답으로 벌금선 아래가 되어 EXIT & 275 \\
K6 & 침묵 & Astra 4103 R1 & 세 번 틀린 판에서 보인 메시지는 이것 하나였고 TALK 응답 61번 중 54번이 DONE뿐 & 32 \\
P1 & 약속 & Opus 4104 R1 & 가장 적게 가진 자리가 1라운드에 서로 메워 주자고 제안했고 넷 모두 끝까지 남음 & 1208 \\
P2 & 셈 & Opus 4103 R6 & 요청과 이전이 같은 틱에 쓰였다며 상대 잔량을 셈하고 더 보내지 않음 & 1084 \\
P5 & 약속 & Sol 4104 R2 & 제삼의 자리가 반환 요청을 거듦; 반환은 같은 틱에 쓰임 & 1370 \\
P6 & 약속 & Astra 4102 R6 & 1라운드의 약속을 제삼의 자리가 6라운드에 꺼냄; 다음 틱에 두 자리가 6,000씩 보냄 & 490 \\
P8 & 조율 & Fable 4102 R7 & solver에게 먼저 주고 다른 자리의 보충은 오답일 때로 미룸; 7, 8라운드 정답, 넷 모두 끝까지 남음 & 1048 \\
P10 & 조율 & Astra 4101 R2 & 한 자리가 자기는 agent-17을 채우고 agent-6에게 agent-11을 맡김; 다음 틱에 4,000이 감 & 214 \\
I1 & 학습 & Sol 4101 R3 & 앞 라운드의 숨은 규칙에서 순서 가설을 세웠고 4라운드에 반례를 보고 거둠 & 125 \\
I4 & 풀이 버퍼 & Opus 4104 R2 & solver에게 8,000씩 두 번 주었으나 풀이는 1,766을 쓰고 틀림 & 1296 \\
I5 & 학습 & Sol 4103 R8 & 떠나는 자리가 다음 solver에게 볼 점을 남김; 그 solver는 틀림 & 899 \\
H1 & 속말 & Fable 4101 R1 & 공용 방에는 넷이 함께 끝나기를 적고 같은 응답의 EXIT: NO 이유에는 끝 할당을 적음 & 56 \\
H3 & 받고 이탈 & Opus 4101 R2 & 3,000을 주며 떠나도 된다고 덧붙임; 받은 자리는 다음 라운드에 EXIT(E8) & 128 \\
H4 & 이탈 & Opus 4101 R6 & 마지막 두 자리가 같은 TALK에서 모두 EXIT를 씀; solver의 EXIT는 막혀 다음 라운드에 떠남(E4) & 259 \\
O1 & 같은 말 & Fable 4105 R4 & 오답 없는 판에서 공용 방 메시지의 80\%가 다른 자리의 메시지와 이름만 다르고 같음; 4라운드에는 두 자리가 같은 틱에 같은 문장을 씀 & 152 \\
\bottomrule
\multicolumn{5}{@{}p{\textwidth}@{}}{\footnotesize 표에 없는 번호(S2, S3, S5, S8, S9, S10, A3, A6, E1, E7, E9, D4, C2, C8, S12, S14, K2, K4, K5, P3, P4, P7, P9, I2, I3, H2, H5, H6, O2)는 같은 장면의 다른 인용이다.} \\
\end{tabularx}
